% !TeX program = lualatex
% !TeX encoding = utf8
% !TeX spellcheck = uk_UA
% !TeX root = ../main.tex

%=========================================================
\chapter{Магнітне поле у вакуумі}\label{\currfilebase}\hypertarget{MagFieldVacuum}{}
\graphicspath{{\currfilebase/Pictures}}
%=========================================================
\epigraph{Магнітне поле не діє на нерухомий заряд і не виконує роботи над рухомим --- і водночас саме воно обертає рамки двигунів і тримає
	стрілку компаса. Ця позірна суперечність зникає, щойно зрозуміти: магнітне поле не додає енергії системі, а лише \enquote{перенаправляє} рух.
}{Переосмислення ідей сучасної фізики}

%% --------------------------------------------------------
\section{Означення}\hypertarget{magnetic-field-definitions}{}
%% --------------------------------------------------------

\begin{wrapstuff}[type=figure, o, top=2, width=0.4\linewidth,]
    \includegraphics[width=\linewidth]{Oersted_experiment}
    \caption{Установка Ерстеда}\label{pic:Oersted_experimente}
\end{wrapstuff}
У 1820~р. датський фізик Г.~Х.~Ерстед виявив, що магнітна стрілка компаса, розташована поблизу провідника, відхиляється, щойно по
провіднику пускають електричний струм (рис.~\ref{pic:Oersted_experimente}). Цей дослід став першим експериментальним доказом того, що електричний струм впливає на магніт, тобто що
рухомі заряди й магнетизм --- взаємопов'язані явища.

До досліду Ерстеда електрику й магнетизм вважали двома окремими, ніяк не пов'язаними розділами фізики. Відкриття Ерстеда --- відправна точка об'єднання цих двох явищ в єдину теорію
електромагнетизму, вершиною якого стануть згодом рівняння Максвелла.

\emphx{Магнітним полем} називають силове поле, що діє на \emphx{рухомі} заряди (а отже, і на електричні струми, і на тіла, які мають магнітний
момент). Ключова відмінність від електричного поля --- магнітне поле не діє на \emphx{нерухомі} заряди: сила, з якою воно діє на заряд, з'являється
лише тоді, коли заряд рухається (рис.~\ref{pic:ChargeFields}).

\begin{SCfigure}[1.5][h!]\centering
	\localinput{chargeFields.tikz}
	\caption{Заряд, що рухається зі швидкістю $\vect v$, є джерелом як електричного, так і магнітного
		поля\label{pic:ChargeFields}}
\end{SCfigure}

Магнітне поле створюється рухомими зарядами, тобто електричним струмом; що саме \emphx{рух} заряду, а не наявність провідника, породжує поле, показують досліди Роуленда й Ейхенвальда (с.~\pageref{hist:RowlandEichenwald}). Якщо струми не змінюються в часі, вони створюють \emphx{постійне} (стале)
магнітне поле --- саме такі поля ми і вивчатимемо в цьому розділі, залишаючи змінні в часі поля для подальших розділів курсу.

%% --------------------------------------------------------
\section{Характеристика магнітного поля}\hypertarget{magnetic-field-characteristics}{}
%% --------------------------------------------------------

На відміну від електростатики, де характеристику поля --- напруженість $\Efield$ --- зручно ввести через силу, що діє на \emphx{пробний заряд}
($\Efield = \vect F/q$), для магнітного поля такий підхід не спрацьовує безпосередньо: магнітних зарядів (\emphx{магнітних монополів}) у природі досі не
виявлено, хоча їх шукають.

Натомість у природі існують \emphx{магнітні диполі} --- магнітна стрілка, коловий виток зі струмом тощо. Скориставшись аналогією з моментом сил,
що діє на електричний диполь в електричному полі, $\vect M = \left[\vect p_e\times\Efield\right]$, вводять характеристику магнітного поля --- вектор
$\Bfield$ --- через момент сил, що діє на магнітний диполь із моментом $\vect p_m$:
\begin{equation}\label{eq:TorqueDipoleAnalogy}
	\vect M = \left[\vect p_m\times\Bfield\right], \qquad M_{\max} = p_m B.
\end{equation}
Отже,
\begin{equation*}
	B = \frac{M_{\max}}{p_m}:
\end{equation*}
величина вектора індукції чисельно дорівнює максимальному обертальному моменту, що діє на одиничний магнітний момент, вміщений у це поле. Напрям $\Bfield$ збігається з напрямом, уздовж якого встановлюється магнітний момент у положенні стійкої рівноваги (див.~\eqref{eq:PotentialEnergyDipole}).

Точне означення вектора $\vect p_m$ дамо нижче (§~\ref{sec:MagMoment}); поки що досить сприймати його як міру \enquote{інтенсивності} магнітного
диполя --- так само як електричний дипольний момент $\vect p_e=\sum_i q_i\vect r_i$ \eqref{eq:el_momentum} характеризує систему зарядів, не будучи
при цьому їхнім сумарним зарядом.

\begin{Historical}
	Вектор $\Bfield$ називають не \emphx{напруженістю}, а \emphx{індукцією} магнітного поля --- на відміну від $\Efield$. Причина --- історична. Напруженістю магнітного поля спершу називали вектор $\Hfield$, який вводили через силу, що діє на магнітний полюс (за аналогією з $\Efield$ і законом Кулона для полюсів). Термін \enquote{індукція} закріпився за вектором $\Bfield$, бо в речовині це поле \emphx{індукується} намагнічуванням: $\Bfield$ відрізняється від $\Hfield$ на внесок намагніченості (гл.~\ref{MagFieldMedia}). У вакуумі в гауссовій системі $\Bfield$ і $\Hfield$ збігаються, але саме $\Bfield$ входить у силу Лоренца, тому саме він є фундаментальною характеристикою поля.
\end{Historical}

У гауссовій системі одиниць величину магнітного поля вимірюють у \emphx{гаусах} (Гс), а в СІ --- у \emphx{теслах} (Тл):
\begin{equation*}
	1\ \text{Тл} = 10^4\ \text{Гс}.
\end{equation*}
Для орієнтира: магнітне поле Землі на поверхні --- близько $0{,}5$~Гс, поле звичайного постійного магніту --- сотні Гс, поле надпровідного магніту
в томографі --- одиниці Тл.

%% --------------------------------------------------------
\section{Дія магнітного поля на заряджені частинки та струми} \hypertarget{magnetic-force-on-charges-currents}{}
%% --------------------------------------------------------

\subsection*{Сила Лоренца та сила Ампера}

\emphx{Магнітною складовою сили Лоренца} називають силу, що діє на рухомий заряд $q$ з боку магнітного поля:
\begin{equation}\label{eq:LorentzMagnetic}
		\vect F = q\left[\frac{\vect v}{c}\times\Bfield\right].
\end{equation}
Повна сила, що діє на заряд у присутності як електричного, так і магнітного поля (власне і є \emphx{сила Лоренца}), --- це сума електричної та
магнітної складових:
\begin{equation*}
	\tcbhighmath{
		\vect F = q\left(\Efield + \left[\frac{\vect v}{c}\times\Bfield\right]\right).
	}
\end{equation*}

\emphx{Силою Ампера} називають силу, що діє з боку магнітного поля не на окремий заряд, а на струм. Якщо в елементі об'єму $dV$ провідника
густина струму дорівнює $\vect j$, то, підсумовуючи внесок~\eqref{eq:LorentzMagnetic} від усіх носіїв заряду в цьому об'ємі ($dq\,\vect v \to
\vect j\, dV$, див. нижче), маємо
\begin{equation}\label{eq:AmpereForceVol}
	\tcbhighmath{
		d\vect F = \frac1c\left[\vect j\, dV\times\Bfield\right],
	}
\end{equation}
де величину $\vect j\, dV$ називають \emphx{елементом об'ємного струму}.

%% --------------------------------------------------------
\subsection*{Зв'язок сили Лоренца та сили Ампера}
%% --------------------------------------------------------


Покажемо явно, що формула~\eqref{eq:AmpereForceVol} --- прямий наслідок магнітної сили Лоренца~\eqref{eq:LorentzMagnetic}. Сила Лоренца, що діє на
заряд $dq$, який рухається зі швидкістю $\vect v$, дорівнює $d\vect F = \left[dq\,\vect v/c \times\Bfield\right]$. Оскільки $dq\,\vect v = \rho\vect
v\, dV = \vect j\, dV$ (гл.~\ref{SteadyCurrent}; тут $\rho$ --- густина заряду саме \emphx{рухомих} носіїв: у нейтральному провіднику повна густина заряду дорівнює нулю, а струм є), звідси одразу отримуємо
формулу~\eqref{eq:AmpereForceVol}. Для \emphx{окремого} рухомого заряду $q$ роль \enquote{елемента струму} відіграє добуток $q\vect v$.

%%% --------------------------------------------------------
\subsection*{Елемент струму}
%%% --------------------------------------------------------


Якщо в задачі не цікавляться внутрішньою будовою провідника й розподілом струму в його товщі, елемент об'ємного струму зручно спростити до
\emphx{елемента лінійного струму}. Нехай струм тече тонким провідником із площею поперечного перерізу $S$ (рис.~\ref{pic:CurrentElement}). Уведемо
вектор ділянки провідника завдовжки $d\ell$ за формулою $d\vect\ell = \vect n\, d\ell$, де $\vect n$ --- одиничний вектор уздовж осі провідника.
Тоді $\vect j = j\vect n$, а $I = jS$, тому
\begin{equation}\label{eq:LinearCurrentElement}
	\vect j\, dV = j\vect n\, S\, d\ell = I\, d\vect\ell,
\end{equation}
і сила Ампера для елемента лінійного струму
\begin{equation}\label{eq:AmpereForceLin}
	\tcbhighmath{
		d\vect F = \frac1c\left[I\, d\vect\ell\times\Bfield\right].
	}
\end{equation}

\begin{SCfigure}[3.5][h!]\centering
	\localinput{CurrentElementSpy.tikz}
	\caption{Елемент лінійного струму $I\, d\vect\ell$: збільшений фрагмент показує напрям
		$\vect j$ уздовж осі провідника перерізу $S$\label{pic:CurrentElement}}
\end{SCfigure}

Залежно від того, як саме розподілений струм, зручно користуватись однією з трьох еквівалентних форм елемента струму
(рис.~\ref{pic:CurrentElementTypes}): \emphx{об'ємним} струмом $\vect j\, dV$ (коли важливий розподіл струму по товщі провідника), \emphx{лінійним}
струмом $I\, d\vect\ell$ (коли провідник тонкий, і досить знати лише повний струм у ньому) або \emphx{поверхневим} струмом $\vect i\, dS$ (коли
струм тече тонким шаром по поверхні --- наприклад, струми в надпровідниках чи високочастотні струми в провідниках); тут $i = I/\ell_{\perp}$ --- поверхнева густина струму (сила
струму, що припадає на одиницю довжини лінії, перпендикулярної до напряму струму).

\begin{figure}[h!]
	\centering
	\begin{minipage}[t]{0.32\linewidth}\centering
		\localinput{VolCurrentElement.tikz}\\
		\footnotesize Об'ємний струм $\vect j\, dV$
	\end{minipage}\hfill
	\begin{minipage}[t]{0.32\linewidth}\centering
		\localinput{LinCurrentElement.tikz}\\
		\footnotesize Лінійний струм $I\, d\vect\ell$
	\end{minipage}\hfill
	\begin{minipage}[t]{0.32\linewidth}\centering
		\localinput{SurfCurrentElement.tikz}\\
		\footnotesize Поверхневий струм $\vect i\, dS$
	\end{minipage}
	\caption{Три еквівалентні форми запису елемента струму\label{pic:CurrentElementTypes}}
\end{figure}


%%% --------------------------------------------------------
\subsection*{Робота магнітного поля}
%%% --------------------------------------------------------


Обчислимо роботу магнітної складової сили Лоренца над рухомим зарядом:
\begin{equation*}
	\delta A = \vect F\cdot d\vect r = \vect F\cdot \vect v\, dt = \left[\frac{q\vect v}{c}\times\Bfield\right]\cdot \vect v\, dt.
\end{equation*}
Скориставшись властивістю мішаного добутку, $\vect a\cdot[\vect b\times\vect w] = \vect w\cdot[\vect a\times\vect b]$, маємо $\vect
v\cdot\left[\vect v\times\Bfield\right] = \Bfield\cdot\left[\vect v\times\vect v\right] = 0$ (оскільки $\left[\vect v\times\vect v\right]\equiv
0$), тому
\begin{equation}\label{eq:WorkMagnetic}
	\tcbhighmath{
		\delta A = 0.
	}
\end{equation}
За теоремою про зміну кінетичної енергії, $A = \Delta\left(mv^2/2\right) = 0$, тобто \emphx{кінетична енергія зарядженої частинки в магнітному
	полі не змінюється}: магнітне поле не виконує роботи над частинкою (рис.~\ref{pic:CyclotronicTrajectory} --- траєкторія зарядженої частинки в
однорідному магнітному полі: сила $\vect F$ завжди перпендикулярна до швидкості $\vect v$ і лише викривляє траєкторію, не змінюючи модуля
швидкості).

\begin{SCfigure}[3.5][h!]\centering
	\localinput{CycletronicTrajectory.tikz}
	\caption{Кругова траєкторія зарядженої частинки в однорідному магнітному полі: сила
		завжди перпендикулярна до швидкості\label{pic:CyclotronicTrajectory}}
\end{SCfigure}

%% --------------------------------------------------------
\section{Закон Біо--Савара--Лапласа}\hypertarget{biot-savart-law}{}
%% --------------------------------------------------------

\subsection*{Формулювання закону}

Закон Біо-Савара (встановлений у 1820~р. Ж.-Б.~Біо і Ф.~Саваром на основі аналізу експериментальних даних, а математично узагальнений
П.-С.~Лапласом) визначає магнітне поле, що створюється елементом струму. Якщо радіус-вектор точки спостереження $P$ відносно розглянутого елемента
струму дорівнює $\vect{\mathcal r}$ (рис.~\ref{pic:BSL}), то поле, створюване елементом об'ємного струму $\vect j\, dV'$, дорівнює
\begin{equation}\label{eq:BSL}
	\tcbhighmath{
		d\Bfield = \frac1c\, \frac{\left[\vect j\, dV'\times\vect{\mathcal r}\right]}{\mathcal r^3}.
	}
\end{equation}

\begin{SCfigure}[1.5][h!]\centering
	\localinput{BSL.tikz}
	\caption{Поле $d\Bfield$, створюване елементом струму $\vect j\, dV'$, за законом
		Біо--Савара--Лапласа\label{pic:BSL}}
\end{SCfigure}

Як і електричне поле, магнітне поле підкоряється \emphx{принципу суперпозиції}: поле довільної системи струмів --- це векторна сума (інтеграл)
внесків від усіх елементів струму:
\begin{equation}\label{eq:Superposition}
	\Bfield = \int d\Bfield.
\end{equation}

\begin{Attention}
	Закон Біо--Савара--Лапласа --- це магнітний аналог закону Кулона: обидва є емпіричними законами, обидва підкоряються принципу суперпозиції, і в обох поле спадає з відстанню як $1/\mathcal r^2$. Головна відмінність --- геометрична: поле точкового заряду напрямлене вздовж $\vect{\mathcal r}$ (радіально), тоді як поле елемента струму, через векторний добуток, перпендикулярне і до струму, і до $\vect{\mathcal r}$ (тобто \enquote{закручується} навколо провідника). Є й фізична відмінність: точковий заряд існує сам по собі, а ізольований елемент струму --- ні (стаціонарні струми замкнені, $\divg\vect j = 0$), тому закон~\eqref{eq:BSL} можна перевірити на досліді лише для замкнених контурів.
\end{Attention}

%%% --------------------------------------------------------
\subsection*{Відносність величини магнітного поля}
%%% --------------------------------------------------------


Розглянемо один і той самий рухомий заряд з погляду двох різних спостерігачів: спостерігача $K'$, що рухається \emphx{разом} із зарядом, і
спостерігача $K$, відносно якого заряд рухається зі швидкістю $\vect v$. Відносно спостерігача $K'$ заряд нерухомий, тому магнітного поля навколо
нього немає (є лише електричне поле $\Efield$). Відносно ж спостерігача $K$, за законом Біо--Савара--Лапласа для точкового заряду ($\vect j\, dV
\to q\vect v$),
\begin{equation}\label{eq:BFromBprime}
	\Bfield = \frac1c\, \frac{q\vect v\times\vect{\mathcal r}}{\mathcal r^3} = \frac{\vect v}{c}\times \frac{q\vect{\mathcal r}}{\mathcal r^3} =
	\frac{\vect v}{c}\times\Efield.
\end{equation}

\begin{Attention}
	Той самий заряд для одного спостерігача створює \emphx{лише} електричне поле, а для іншого --- \emphx{і} електричне, \emphx{і} магнітне.
	Оскільки саме поле --- фізично реальний об'єкт, а не розрахункова абстракція, цей результат означає, що поділ електромагнітного поля на
	\enquote{електричну} і \enquote{магнітну} складові \emphx{залежить від системи відліку}: електричне й магнітне поля --- не два окремі поля, а
	дві грані одного й того самого \emphx{електромагнітного поля}. Строге і повне розкриття цього факту дає спеціальна теорія відносності, де
	$\Efield$ і $\Bfield$ об'єднуються в  єдиний тензор електромагнітного поля; формула~\eqref{eq:BFromBprime} --- це лише перше (нерелятивістське,
	$v\ll c$) наближення до точного релятивістського закону перетворення полів (див.~розділ \enquote{\nameref{Relativity}}).
\end{Attention}

\begin{Historical}\label{hist:RowlandEichenwald}
	\emphx{Досліди Роуленда й Ейхенвальда.} Чи створює магнітне поле сам рух зарядженого тіла, а не лише струм у провіднику? Від відповіді залежить справедливість формули~\eqref{eq:BFromBprime}. Її дали досліди Г.~Роуленда (1876~р., у лабораторії Г.~Гельмгольца в Берліні) та О.~Ейхенвальда (1901--1903~рр.), в яких спостерігали магнітну дію \emphx{конвекційного струму}, тобто заряджених тіл, що рухаються: заряджений диск, який швидко обертається, діє на магнітну стрілку так само як і колові струми (рис.~\ref{pic:EichenwaldExperiment}).

	Оцінимо ефект. Нехай диск радіуса $R$ має поверхневу густину заряду $\sigma$ і обертається з кутовою швидкістю $\omega$. Кільце радіуса $r$ і завширшки $dr$ дає струм $dI = \sigma\omega r\, dr$, а його поле в центрі, за законом Біо--Савара--Лапласа, дорівнює $dB = 2\pi\, dI/(cr)$. Звідси
	\begin{equation*}
		B = \frac{2\pi\sigma\omega R}{c} = \frac{2Q\omega}{cR}.
	\end{equation*}
	Для реальних параметрів це поле надзвичайно мале (задача~\ref{prob:RowlandDisk}), тому потрібні були чутливі астатичні магнітометри й ретельне екранування від земного поля. Дослід стояв на межі чутливості приладів і не раз піддавався сумніву (на початку ХХ~ст. В.~Кремьє спершу не зміг його відтворити), проте Ейхенвальд суттєво підвищив точність і підтвердив, що магнітна дія конвекційного струму збігається з дією рівного йому струму провідності. 
	% Він також дослідив магнітну дію діелектрика, що рухається в електричному полі (\emphx{струм Рентгена}): магнітне поле створює й рух поляризованого діелектрика, тобто зв'язаних зарядів. Це не те саме, що \emphx{струм зміщення}, який не пов'язаний з рухом зарядів узагалі (гл.~\ref{MaxwellEqns}).

	% \begin{center}
	% 	% TODO: ЗАГЛУШКА, замінити на \localinput{EichenwaldExperiment.tikz}
	% 	\begin{tikzpicture}
	% 		\node[draw, dashed, gray, minimum width=0.7\linewidth, minimum height=3.5cm, align=center, text=gray]
	% 		{Заглушка для рисунка:\\ заряджений диск, що обертається,\\ і магнітна стрілка (астатична система)};
	% 	\end{tikzpicture}
	% 	\captionof{figure}{Схема дослідів Роуленда й Ейхенвальда: заряджений диск, що обертається, діє на магнітну стрілку\label{pic:EichenwaldExperiment}}
	% \end{center}
\end{Historical}

\begin{Example}{Поле нескінченного прямого провідника.}
Знайдемо магнітне поле на відстані $r$ від нескінченно довгого прямолінійного провідника
зі струмом $I$. Скориставшись результатом задачі про поле зарядженої нитки (гл.~\ref{SteadyEfield}) як геометричною підказкою, --- інтеграл
за формою той самий, --- виконаємо пряме інтегрування закону Біо--Савара--Лапласа вздовж провідника.

\begin{center}
    \localinput{long-wire.tikz}
    \captionof{figure}{До обчислення магнітного поля нескінченного провідника зі струмом.}
    \label{tikz:long-wire}
\end{center}

Нехай точка спостереження перебуває на
відстані $r$ від провідника; параметризуємо елемент струму кутом $\alpha$ між $\vect{\mathcal r}$ і перпендикуляром, опущеним із точки
спостереження на провідник, $d\ell = r\, d\alpha/\cos^2\alpha$, $\mathcal r = r/\cos\alpha$; кожен елемент дає внесок, напрямлений однаково
(перпендикулярно до площини, що містить провідник і точку спостереження), тому поля просто додаються:
\begin{equation*}
	B = \int\limits_{-\pi/2}^{\pi/2} \frac{I\, d\ell\, \sin(90^\circ-\alpha)}{c\,\mathcal r^2} = \frac{I}{cr}\int\limits_{-\pi/2}^{\pi/2}
	\cos\alpha\, d\alpha = \frac{2I}{cr}.
\end{equation*}
Отже,
\begin{equation}\label{eq:InfiniteWireResult}
		B = \frac{2I}{cr}.
\end{equation}

\end{Example}

\begin{Example}{Поле відрізка провідника}
Визначимо магнітне поле в точці $P$ на відстані $r$ від \emphx{скінченного} відрізка провідника зі струмом;
положення точки визначається кутами $\alpha_1$ і $\alpha_2$ між перпендикуляром, опущеним з $P$ на провідник, і напрямами на кінці відрізка.
%(рис.~\ref{pic:FiniteWireAngles}).
Обчислення повторює попередній приклад, тільки тепер межі інтегрування скінченні:
\begin{equation*}
		B = \frac{I}{cr}\left(\sin\alpha_2 - \sin\left(-\alpha_1\right)\right) = \frac{I}{cr}\left(\sin\alpha_1+\sin\alpha_2\right).
\end{equation*}
При $\alpha_1,\alpha_2 \to \pi/2$ (провідник стає нескінченним в обидва боки) ця формула, як і мало бути, переходить у результат~\eqref{eq:InfiniteWireResult}.
\begin{center}
	\localinput{finite-wire-angles.tikz}
%	\captionof{figure}{Геометрія задачі про поле скінченного відрізка провідника\label{pic:FiniteWireAngles}}
\end{center}
\end{Example}

%%% --------------------------------------------------------
\subsection*{Досліди Ампера. Взаємодія струмів}
%%% --------------------------------------------------------


У 1820--1826~рр. А.~Ампер вивчав силову взаємодію струмів і сформулював закон взаємодії елементів струму, з якого випливає й вираз для сили, що діє на елемент струму в магнітному полі. Оскільки відокремлений елемент струму
експериментально створити не можна, Ампер вивчав вплив паралельних провідників один на одного та поведінку замкнених контурів різної форми в
магнітному полі. Він встановив дослідним шляхом, що паралельні провідники зі струмами, що течуть в \emphx{одному} напрямку, \emphx{притягуються},
а в \emphx{протилежних} напрямках --- \emphx{відштовхуються}.
%  (рис.~\ref{pic:AmpereForce}).

% \begin{wrapstuff}[type=figure, o, top=2, width=0.4\linewidth]
% 	\centering
% 	\includegraphics[width=\linewidth]{AmpereForce.pdf}
% 	\caption{Сила взаємодії паралельних провідників}
% 	\label{pic:AmpereForce}
% \end{wrapstuff}

Скориставшись результатом~\eqref{eq:InfiniteWireResult} і законом Ампера~\eqref{eq:AmpereForceLin}, легко отримати кількісний вираз для сили
взаємодії двох нескінченних паралельних провідників зі струмами $I_1$ і $I_2$ на відстані $d$ один від одного, у розрахунку на одиницю довжини:
\begin{equation}\label{eq:AmpereParallelWires}
		\frac{F}{\ell} = \frac{2I_1I_2}{c^2 d}.
\end{equation}

\begin{Historical}
	Саме формула~\eqref{eq:AmpereParallelWires} історично лежала в основі означення одиниці сили струму --- \emphx{Ампера} --- в старій версії
	Міжнародної системи одиниць (до реформи 2019~р.): 1~А означав силу струму, за якої два безмежно довгі паралельні провідники мізерно малого
	перерізу, розташовані у вакуумі на відстані 1~м один від одного, взаємодіють із силою $2\cdot10^{-7}$~Н на кожен метр довжини.
\end{Historical}

%% --------------------------------------------------------
\section{Вектор-потенціал магнітного поля}\hypertarget{magnetic-vector-potential}{}
%% --------------------------------------------------------

Запишемо закон Біо--Савара--Лапласа~\eqref{eq:BSL}--\eqref{eq:Superposition} для довільного розподілу струму $\vect j(\vect r')$
(рис.~\ref{pic:FieldPoint}), користуючись позначенням $\vect{\mathcal r} = \vect r - \vect r'$:
\begin{equation*}
	\Bfield(\vect r) = \iiint\limits_{V'} \frac1c\, \frac{\vect j(\vect r')\, dV' \times(\vect r-\vect r')}{|\vect r - \vect r'|^3}.
\end{equation*}

\begin{SCfigure}[1][h!]\centering
	\localinput{FieldPoint.tikz}
	\caption{Позначення для точки спостереження $\vect r$ і точки джерела $\vect r'$\label{pic:FieldPoint}}
\end{SCfigure}

Скористаємось тотожністю $\dfrac{\vect r - \vect r'}{|\vect r-\vect r'|^3} = -\grad_{\vect r} \dfrac1{|\vect r - \vect r'|}$ (тут $\grad$ діє на
координати точки спостереження $\vect r$), а також тотожністю $\rot(\phi\vect C) = \phi\rot\vect C + \left[\grad\phi\times\vect C\right]$.
Оскільки $\vect j(\vect r')$ не залежить від $\vect r$, $\rot_{\vect r}\vect j(\vect r') = 0$, і
\begin{equation*}
	\Bfield(\vect r) = \int\limits_{V'} \frac1c \grad\frac1{|\vect r-\vect r'|}\times\vect j(\vect r')\, dV' = \rot\frac1c\int\limits_{V'}
	\frac{\vect j(\vect r')\, dV'}{|\vect r-\vect r'|}.
\end{equation*}
Отже, магнітне поле можна подати у вигляді ротора деякого допоміжного векторного поля:
\begin{equation}\label{eq:BrotA}
	\tcbhighmath{
		\Bfield = \rot\vect A,
	}
\end{equation}
де $\vect A$ називають \emphx{вектор-потенціалом} магнітного поля:
\begin{equation}\label{eq:VectorPotentialDef}
	\tcbhighmath{
		\vect A(\vect r) = \frac1c\int\limits_{V'} \frac{\vect j(\vect r')\, dV'}{|\vect r-\vect r'|}, \qquad \vect A(\vect r) = \frac1c
		\sum\limits_i \frac{q_i\vect v_i}{|\vect r - \vect r_i|}
	}
\end{equation}
(друга формула --- для системи точкових зарядів).

\begin{Attention}
	Вектор-потенціал $\vect A$ відіграє в магнітостатиці ту саму роль, що й скалярний потенціал $\phi$ в електростатиці: обидва --- допоміжні
	величини, через які поле виражається простіше, ніж напряму. Принципова відмінність --- $\vect A$ векторний, а $\phi$ скалярний, що
	відображає \emphx{вихровий} характер магнітного поля на противагу \emphx{потенціальному} характеру електростатичного.
\end{Attention}

\subsection*{Калібрування вектор-потенціалу}

Вектор-потенціал вводиться \emphx{неоднозначно}: два вектор-потенціали $\vect A$ і $\vect A' = \vect A + \grad f(\vect r)$ (де $f$ --- будь-яка
скалярна функція) призводять до \emphx{одного й того самого} магнітного поля $\Bfield$, оскільки $\rot(\grad f) \equiv 0$ для будь-якої функції
$f$. Цю свободу вибору (\emphx{калібрувальну свободу}) можна використати, щоб накласти на $\vect A$ якесь додаткове зручне обмеження. Найзручніше в
магнітостатиці --- накласти умову
\begin{equation}\label{eq:CoulombGauge}
	\tcbhighmath{
		\divg\vect A = 0
	}
\end{equation}
--- її називають \emphx{кулонівським калібруванням}. Саме таке калібрування задовольняє вираз~\eqref{eq:VectorPotentialDef}, отриманий напряму
із закону Біо--Савара--Лапласа (у цьому неважко переконатись, узявши дивергенцію і скориставшись рівнянням неперервності для стаціонарного
струму, $\divg\vect j = 0$; струми вважаємо локалізованими в скінченній області).

%% --------------------------------------------------------
\section{Теореми магнітостатики}\hypertarget{magnetostatics-theorems}{}
%% --------------------------------------------------------

\subsection*{Теорема Гаусса для магнітного поля}

Маючи на увазі тотожність $\divg\rot\vect A \equiv 0$ (дивергенція будь-якого ротора тотожно дорівнює нулю) і формулу $\Bfield = \rot\vect A$,
відразу отримуємо \emphx{теорему Гаусса для магнітного поля в диференціальній формі}:
\begin{equation}\label{eq:GaussMagDiff}
	\tcbhighmath{
		\divg\Bfield = 0.
	}
\end{equation}
Застосовуючи теорему Остроградського--Гаусса, отримуємо інтегральну форму:
\begin{equation*}
	\tcbhighmath{
		\oiint\limits_S \Bfield\cdot d\vect S = 0.
	}
\end{equation*}

\begin{Attention}
	Теорема Гаусса для $\Bfield$ --- це математичний вираз того самого факту, з якого ми почали розділ: \emphx{вільних магнітних зарядів немає} (у жодному досліді їх досі не виявлено).
	Силові лінії індукції магнітного поля не мають ні витоків, ні стоків --- вони або замкнені самі на себе, або йдуть з нескінченності в
	нескінченність, але ніколи не починаються й не закінчуються в певній точці простору (на відміну від ліній $\Efield$, які починаються й
	закінчуються на електричних зарядах).
\end{Attention}

\subsection*{Теорема про циркуляцію: диференціальна форма}

Знайдемо ротор вектора $\Bfield$, скориставшись відомою тотожністю $\rot\rot\vect A = \grad(\divg\vect A) - \nabla^2\vect A$ і кулонівським
калібруванням~\eqref{eq:CoulombGauge} ($\divg\vect A = 0$):
\begin{equation}\label{eq:CurlBDiff}
	\rot\Bfield = \rot\rot\vect A = \grad(\divg\vect A) - \nabla^2\vect A = -\nabla^2\vect A.
\end{equation}
З іншого боку, можна показати (за прямою аналогією з електростатикою), що кожна декартова компонента вектора $\vect A$ \eqref{eq:VectorPotentialDef} задовольняє рівняння Пуассона з джерелом $\vect j/c$:
\begin{equation*}
	\nabla^2\phi = -4\pi\rho, \quad \phi = \iiint\limits_{V'} \frac{\rho\, dV'}{|\vect r-\vect r'|} \qquad\Longleftrightarrow\qquad \nabla^2\vect A
	= -\frac{4\pi}c\vect j, \quad \vect A = \frac1c \iiint\limits_{V'} \frac{\vect j\, dV'}{|\vect r-\vect r'|}.
\end{equation*}

\begin{Attention}
	\emphx{Однакові за формою рівняння за однакових граничних умов мають однакові розв'язки} --- це загальний і дуже потужний принцип математичної фізики: рівняння
	Пуассона для кожної компоненти $\vect A$ математично \emphx{тотожне} рівнянню Пуассона для потенціалу $\phi$, з тією лише заміною $\rho \to
	j_x/c,\ j_y/c,\ j_z/c$. Тому математичний апарат, розвинений для електростатики (методи розв'язання рівняння Пуассона тощо), переноситься на задачі магнітостатики --- варто лише правильно \enquote{перекласти} одну задачу в термінах іншої і врахувати, що граничні умови для $\vect A$ на межах середовищ інші, ніж для $\phi$.
\end{Attention}

Підставляючи $\nabla^2\vect A = -4\pi\vect j/c$ у~\eqref{eq:CurlBDiff}, отримуємо \emphx{теорему про циркуляцію магнітного поля в диференціальній
	формі}:
\begin{equation}\label{eq:CirculationDiff}
	\tcbhighmath{
		\rot\Bfield = \frac{4\pi}c\vect j.
	}
\end{equation}
Зауважимо, що цю теорему виведено для \emphx{стаціонарних} струмів ($\divg\vect j = 0$). Для змінних у часі полів її доведеться доповнити \emphx{струмом зміщення} (гл.~\ref{MaxwellEqns}).

\subsection*{Теорема про циркуляцію: інтегральна форма}

Застосовуючи до~\eqref{eq:CirculationDiff} теорему Стокса, отримуємо інтегральну форму: циркуляція вектора $\Bfield$ по довільному замкненому
контуру $L$ пропорційна повному струму, що протікає крізь будь-яку поверхню $S$, натягнуту на цей контур
(рис.~\ref{pic:CirculationLoop}):
\begin{equation*}
	\tcbhighmath{
		\oint\limits_L \Bfield\cdot d\vect r = \frac{4\pi}c \iint\limits_S \vect j\cdot d\vect S.
	}
\end{equation*}

\begin{SCfigure}[2][h!]\centering
	\localinput{circulation-loop.tikz}
	\caption{Контур $L$ і натягнута на нього поверхня $S$, крізь яку протікає струм\label{pic:CirculationLoop}}
\end{SCfigure}

У випадку дискретних (лінійних) струмів це співвідношення набуває найпростішого вигляду: циркуляція $\Bfield$ по контуру $L$ пропорційна
\emphx{алгебраїчній} сумі струмів, що охоплюються цим контуром (рис.~\ref{pic:CirculationSigns}):
\begin{equation}\label{eq:CirculationDiscrete}
	\tcbhighmath{
		\oint\limits_L \Bfield\cdot d\vect r = \frac{4\pi}c \sum\limits_i I_i.
	}
\end{equation}

\begin{SCfigure}[1.75][h!]\centering
	\localinput{circulation-signs.tikz}
	\caption{Струм враховують зі знаком \enquote{плюс}, якщо його напрям узгоджений із напрямом
		обходу контуру $L$ за правилом правого гвинта, і зі знаком \enquote{мінус} --- у протилежному випадку\label{pic:CirculationSigns}}
\end{SCfigure}

\begin{Attention}
	Знак кожного струму $I_i$ у сумі~\eqref{eq:CirculationDiscrete} визначають за \emphx{правилом правого гвинта} (правилом буравчика): якщо
	чотири пальці правої руки вказують напрям обходу контуру $L$, то великий палець вказує напрям, у якому струм слід вважати
	\emphx{додатним}. Провідник, що не проходить крізь поверхню, натягнуту на контур (або проходить крізь неї двічі, в обидва боки, як у
	підковоподібному контурі), не робить внеску в циркуляцію: його поле в точках контуру не дорівнює нулю, але циркуляція цього поля по контуру дорівнює нулю.
\end{Attention}


%%% --------------------------------------------------------
\subsection*{Приклади на теорему про циркуляцію}
%%% --------------------------------------------------------


Теорема про циркуляцію, так само як і теорема Гаусса в електростатиці, дає змогу без громіздкого інтегрування знайти поле в задачах з високою
симетрією --- достатньо обрати контур $L$ так, щоб на ньому $B = \const$, а $\Bfield \parallel d\vect r$.

\emphx{Поле нескінченного прямого провідника.} Розглянемо рис.~\ref{tikz:long-wire}. Обираємо контур $L$ --- коло радіуса $r$ у площині, перпендикулярній до провідника, з центром на
провіднику. З міркувань симетрії $\Bfield$ напрямлене по дотичній до цього кола й однакове за модулем усюди на ньому, тому
$\oint_L \Bfield\cdot d\vect r = B\cdot 2\pi r$, а циркуляція охоплює струм $I$; звідси одразу $B = 2I/(cr)$ --- той самий результат, що й ми вже
отримали прямим інтегруванням закону Біо--Савара--Лапласа, формула~\eqref{eq:InfiniteWireResult}, але тепер --- значно простіше.

\emphx{Поле нескінченного соленоїда.} Обираємо прямокутний контур $L$ ($A$, $B$, $C$, $D$), одна сторона якого проходить усередині соленоїда паралельно його осі ($AB$), а
протилежна --- далеко зовні ($CD$), де поле відсутнє (рис.~\ref{pic:solenoidField}); на поперечних сторонах внесок дорівнює нулю (усередині $\Bfield\perp d\vect r$, зовні $\Bfield = 0$). Циркуляція набирається лише з
внутрішньої сторони довжиною $\ell$, $\oint_L \Bfield \cdot d\vect r = B\ell$, а контур охоплює $N$ витків із струмом $I$ кожен, тобто повний струм
$NI$:
\begin{equation*}
	B = \frac{4\pi}{c}\frac{N}{\ell}I = \frac{4\pi}{c}nI,
\end{equation*}
де $n = N/\ell$ --- число витків на одиницю довжини соленоїда. Поле всередині нескінченного соленоїда \emphx{однорідне}, зовні --- дорівнює нулю.

%---------------------------------------------------------
\begin{figure}[h!]\centering
    \includegraphics[width=0.6\linewidth]{solenoidField}
\caption{До розрахунку магнітного поля нескінченно довгого соленоїда.}
\label{pic:solenoidField}
\end{figure}
%---------------------------------------------------------

\emphx{Поле всередині й поза нескінченним циліндричним провідником} (рис.~\ref{pic:CylinderField}), по якому тече струм зі сталою по перерізу
густиною $j$. Усередині (на відстані $r$ від осі, $r<R$) контур $L$ охоплює лише частину струму, $I(r) = j\pi r^2$, тому
\begin{equation*}
	B\cdot 2\pi r = \frac{4\pi}{c}\, j\pi r^2 \quad\Rightarrow\quad B = \frac{2\pi}{c}\,jr \qquad (r < R);
\end{equation*}
поле всередині зростає \emphx{лінійно} з відстанню від осі. Зовні ($r>R$) контур охоплює весь струм $I = j\pi R^2$, тому
\begin{equation*}
	B = \frac{2I}{cr} \qquad (r > R)
\end{equation*}
--- поле спадає так само, як і поле нескінченно тонкого провідника з тим самим повним струмом $I$ (порівняйте з формулою~\eqref{eq:InfiniteWireResult}).

\begin{figure}[h!]\centering
	\localinput{cylinder-field.tikz}
	\caption{Поле нескінченного циліндричного провідника зі струмом: усередині й поза
		провідником\label{pic:CylinderField}}
\end{figure}

%% --------------------------------------------------------
\subsection*{Порівняння законів електро- та магнітостатики у вакуумі}
%% --------------------------------------------------------

Корисно звести головні результати цього й попередніх розділів у порівняльну таблицю: математична структура електро- і магнітостатики у вакуумі
паралельна, хоча фізичний зміст --- принципово різний (табл.~\ref{tab:ElMagCompare}).

\begin{Attention}
	Джерелом електростатичного поля є заряди --- це поле \emphx{потенціальне} (має витоки й стоки, циркуляція завжди нульова). Джерелом
	магнітного поля є струми --- це поле \emphx{вихрове} (не має витоків і стоків, зате має ненульову циркуляцію там, де є струм). Ця структурна
	відмінність --- не випадковість, а наслідок відсутності магнітних зарядів: там, де в електростатиці стояв би заряд $\rho$, у магнітостатиці
	стоїть \emphx{нуль} (теорема Гаусса), а там, де в електростатиці стояв би нуль (потенціальність поля), у магнітостатиці стоїть струм $\vect j$
	(вихровість поля).
\end{Attention}

\begin{table}[h!]\small\centering
	\caption{Порівняння теорем електро- та магнітостатики у вакуумі\label{tab:ElMagCompare}}
	\begin{tblr}{
		colspec={X[l, m]X[c,m]X[c,m]},
        row{odd} = {gray!10},
		row{1} = {c, m, fg=white, bg=themecolorlight, font=\bfseries},
		hlines,
		vlines
		}
		                          & Електростатика                                                & Магнітостатика                                                \\
		Потенціал                & $\Efield = -\grad\phi$ (потенціальне поле)                   & $\Bfield = \rot\vect A$ (вихрове поле)                        \\
		Теорема Гаусса (диф.)    & $\divg\Efield = 4\pi\rho$                                     & $\divg\Bfield = 0$                                            \\
		Теорема Гаусса (інт.)    & $\displaystyle\oiint_S \Efield\cdot d\vect S = 4\pi \!\!\iiint\limits_V\!\! \rho\, dV$ & $\displaystyle\oiint_S \Bfield\cdot d\vect S = 0$             \\
		Теорема про циркуляцію (диф.) & $\rot\Efield = 0$                                        & $\rot\Bfield = \dfrac{4\pi}c\vect j$                          \\
		Теорема про циркуляцію (інт.) & $\displaystyle\oint_L \Efield\cdot d\vect r = 0$          & $\displaystyle\oint_L \Bfield\cdot d\vect r = \dfrac{4\pi}c\!\!\iint\limits_S\!\!\vect j\cdot d\vect S$
	\end{tblr}
\end{table}

%% --------------------------------------------------------
\section{Магнітний момент}\label{sec:MagMoment}\hypertarget{magnetic-moment}{}
%% --------------------------------------------------------

%%% --------------------------------------------------------
\subsection*{Означення}
%%% --------------------------------------------------------


В електростатиці роль \enquote{нульового наближення} відігравав сумарний заряд $q=\sum_i q_i$, а дипольний момент \eqref{eq:el_momentum}
з'являвся вже як поправка --- головна характеристика лише тоді, коли $q=0$. У магнетизмі ситуація радикальніша: магнітних зарядів не існує
взагалі, тож природного аналога величини $q$ немає навіть формально. Найпростіша величина, яка могла б претендувати на цю роль --- сумарна
\enquote{кількість руху заряду} $\sum_i q_i\vect v_i$ --- виявляється (як буде показано нижче, при виведенні~\eqref{eq:VectorPotentialTwoTerms}--\eqref{eq:VectorPotentialFar}) похідною за часом від
електричного дипольного моменту системи
\eqref{eq:el_momentum} і тому в усередненому за часом полі тотожно зникає для будь-якого фінітного руху. Отже, для струмів немає навіть
\enquote{нульового} члена розкладу --- і першою характеристикою, що реально визначає створюване поле, одразу виявляється \emphx{магнітний
момент}: міра того, як заряди рухаються один відносно одного, --- цілком аналогічно до того, як для електричного поля системи з $q=0$ головну
роль відігравало взаємне \emphx{розташування} зарядів \eqref{eq:el_momentum}.

Означають його за прямою аналогією з моментом імпульсу $\vect L = \iiint_V \vect r\times \rho\vect v\, dV$ для розподіленої маси, що рухається:
якщо момент імпульсу --- це \enquote{момент кількості руху}, то магнітний момент --- \enquote{момент кількості струму}:
\begin{equation*}
	\vect L = \iiint\limits_V \vect r\times\rho_m\vect v\, dV, \qquad \vect p_m = \frac1{2c}\iiint\limits_V \vect r\times\rho\vect v\, dV,
\end{equation*}
де $\rho_m$ --- густина маси, а $\rho$ --- густина заряду.

Оскільки густина струму $\vect j = \rho\vect v$, магнітний момент найпростіше записати через об'ємний елемент струму:
\begin{equation}\label{eq:MagMomentVol}
	\tcbhighmath{
		\vect p_m = \frac1{2c}\iiint\limits_V \vect r\times\vect j\, dV.
	}
\end{equation}
Але залежно від того, яким елементом струму (див.~\eqref{eq:LinearCurrentElement}) зручніше описати розподіл струму в конкретній задачі, ця сама
величина набуває різного вигляду:
\begin{align}
	\vect p_m &= \frac{I}{2c}\oint\limits_L \vect r\times d\vect\ell && \text{(лінійний замкнений струм)}, \label{eq:MagMomentLin}\\
	\vect p_m &= \frac1{2c}\sum_i q_i(\vect r_i\times\vect v_i) && \text{(система дискретних точкових зарядів).} \label{eq:MagMomentDiscrete}
\end{align}

%%% --------------------------------------------------------
\subsection*{Магнітний момент колового витка}
%%% --------------------------------------------------------


Обчислимо за формулою~\eqref{eq:MagMomentLin} магнітний момент плоского колового витка радіуса $R$ зі струмом $I$ (рис.~\ref{pic:LoopMagMoment}).
Оскільки $\vect r \perp d\vect\ell$ у кожній точці кола, а $|\vect r\times d\vect\ell| = R\, d\ell$, і напрям цього добутку той самий у кожній
точці (по нормалі $\vect n$ до площини витка), маємо
\begin{equation}\label{eq:MagMomentLoop}
	\tcbhighmath{
		\vect p_m = \frac{I}{2c}\oint\limits_L \vect r\times d\vect\ell = \frac1c I\pi R^2\vect n = \frac1c IS\vect n,
	}
\end{equation}
де $S = \pi R^2$ --- площа витка.

\begin{SCfigure}[2][h!]\centering
	\localinput{loop-magnetic-moment.tikz}
	\caption{До обчислення магнітного моменту колового витка зі струмом\label{pic:LoopMagMoment}}
\end{SCfigure}

\begin{Attention}
	Формула~\eqref{eq:MagMomentLoop} застосовна до витка довільної форми: магнітний момент \emphx{будь-якого} плоского замкненого витка (не обов'язково
	колового) зі струмом $I$ дорівнює $\vect p_m = IS\vect n/c$, де $S$ --- площа, обмежена контуром витка, а напрям нормалі $\vect n$ пов'язаний
	із напрямом обходу контуру струмом за правилом правого гвинта.
\end{Attention}

%%% --------------------------------------------------------
\subsection*{Гіромагнітне відношення}
%%% --------------------------------------------------------


Відношення магнітного моменту зарядженого тіла до його моменту імпульсу називають \emphx{гіромагнітним відношенням}:
\begin{equation*}
	\tcbhighmath{
		\vect p_m = \gamma\vect L.
	}
\end{equation*}
Для будь-якого \emphx{класичного} тіла, у якого відношення заряду до маси однакове в кожному елементі об'єму, гіромагнітне відношення не залежить від закону руху й дорівнює
\begin{equation}\label{eq:GyromagneticClassical}
	\gamma = \frac{Q}{2Mc},
\end{equation}
де $Q$ і $M$ --- повний заряд і маса тіла. Справді, якщо $\rho/\rho_m = Q/M$ у кожній точці, а заряд рухається разом із масою, то
\begin{equation*}
	\vect p_m = \frac1{2c}\iiint\limits_V \vect r\times\rho\vect v\, dV = \frac{Q}{2Mc}\iiint\limits_V \vect r\times\rho_m\vect v\, dV = \frac{Q}{2Mc}\,\vect L.
\end{equation*}

\begin{Historical}
	Класичний результат~\eqref{eq:GyromagneticClassical} чудово підтверджується для макроскопічних тіл, але \emphx{порушується} для електрона:
	його експериментальне гіромагнітне відношення виявляється майже рівно \emphx{вдвічі} більшим за класичне значення $-e/(2m_ec)$, тобто
	$\gamma_e \approx -e/(m_ec)$. Причина --- власний (спіновий) момент імпульсу електрона не пов'язаний із класичним обертанням якогось
	зарядженого тіла навколо осі, а є суто квантовою величиною, для якої \enquote{аномальний} множник 2 ($g$-фактор $g=2$) є наслідком релятивістського квантового рівняння Дірака, а мала поправка ($g = 2{,}0023\ldots$) --- результатом квантової електродинаміки. Вперше на відмінність гіромагнітного відношення від класичного вказали досліди Ейнштейна--де Гааза й Барнетта (1915~р.) з феромагнетиками. Це один із перших історичних натяків на те, що електрон не
	можна уявляти як маленьку заряджену кульку, що обертається.
\end{Historical}

%% --------------------------------------------------------
\subsection*{Обчислення магнітного моменту через момент інерції}
%% --------------------------------------------------------

Зв'язок $\vect p_m = \gamma\vect L$ дає зручний практичний спосіб обчислювати магнітний момент тіл, що обертаються, не вдаючись щоразу до
прямого інтегрування $\vect p_m = \frac1{2c}\iiint_V \vect r\times\vect j\, dV$: досить знайти момент імпульсу $\vect L$ тіла засобами звичайної
механіки твердого тіла. Нагадаємо: для твердого тіла, що обертається з кутовою швидкістю $\vect\omega$ навколо нерухомої осі, момент імпульсу
відносно цієї осі дорівнює $\vect L = J\vect\omega$, де $J$ --- момент інерції тіла відносно осі обертання. Якщо, як і при виведенні
формули~\eqref{eq:GyromagneticClassical}, відношення заряду до маси однакове в кожному елементі об'єму тіла, підставляючи $\vect L=J\vect\omega$
у $\vect p_m=\gamma\vect L$, отримуємо
\begin{equation}\label{eq:MagMomentViaInertia}
		\vect p_m = \frac{Q}{2Mc}\, J\vect\omega.
\end{equation}
Отже, для однорідно (за відношенням $Q/M$) зарядженого тіла, що обертається, обчислення магнітного моменту зводиться до звичної задачі
механіки --- знаходження моменту інерції $J$ відносно осі обертання за відомими формулами (куля, диск, обруч тощо), а не до нового
інтегрування за струмами.

\begin{Example}{Магнітний момент зарядженої кулі, що обертається.}
Знайдемо магнітний момент однорідно зарядженої кулі масою $M$, зарядом $Q$ і радіусом $R$, що обертається з кутовою швидкістю $\omega$ навколо
свого діаметра.

\begin{wrapstuff}[type=figure, width=0.1\linewidth, o]%[6][h!]\centering
	\localinput{gyromagnetic-sphere.tikz}
%	\caption{Однорідно заряджена куля, що обертається навколо власної осі\label{pic:GyromagneticSphere}}
\end{wrapstuff}


 Момент інерції однорідної кулі відносно осі, що проходить через центр, --- відома з
механіки величина $J=\frac25 MR^2$. За формулою~\eqref{eq:MagMomentViaInertia}
\begin{equation*}
	p_m = \frac{Q}{2Mc}\cdot\frac25 MR^2\omega = \frac{QR^2\omega}{5c}.
\end{equation*}
Той самий результат можна отримати прямим інтегруванням $\vect p_m = \frac1{2c}\iiint_V \vect r\times\vect j\, dV$ по об'єму кулі, але шлях
через момент інерції --- значно коротший.
\end{Example}

%% --------------------------------------------------------
\subsection*{Вектор-потенціал на далеких відстанях. Поле магнітного диполя}
%% --------------------------------------------------------

Знайдемо вектор-потенціал системи струмів, зосереджених у малій області простору, на великій відстані від неї ($r \gg r_i$, де $r_i$ ---
характерний розмір системи). Розкладаючи $1/|\vect r - \vect r_i|$ у формулі~\eqref{eq:VectorPotentialDef} в ряд за малим параметром $r_i/r$ і
усереднюючи за часом (для стаціонарного, тобто періодичного або усталеного, руху усереднена похідна будь-якої обмеженої величини дорівнює нулю),
після певних (дещо громіздких, але прямолінійних) векторних перетворень отримуємо головний, дипольний член розкладу.

Нехай вектор-потенціал системи точкових зарядів $q_i$ визначається формулою~\eqref{eq:VectorPotentialDef}:
\begin{equation*}
	\vect A(\vect r) = \frac{1}{c}\sum_i \frac{q_i\vect v_i}{|\vect r - \vect r_i|}.
\end{equation*}
Оскільки заряди рухаються, ця формула дає \emphx{миттєвий} вектор-по\-тен\-ці\-ал, який явно залежить від часу через $\vect r_i(t)$ і $\vect v_i(t)$.
Нас, однак, цікавить не миттєве, а усереднене за часом поле системи --- те, що спостерігається макроскопічно, коли рух зарядів фінітний
(обмежений у просторі) і періодичний або усталений (наприклад, електрони в атомі чи струм у витку). Сенс усереднення в тому, що для
будь-якої \emphx{обмеженої} (фінітної) величини $Q(t)$ її похідна за часом усереднюється до нуля:
\begin{equation}\label{eq:AveragedDerivativeZero}
	\overline{\frac{dQ}{dt}} \equiv \frac{1}{T}\int_0^T \frac{dQ}{dt}\,dt = \frac{Q(T)-Q(0)}{T} \to 0,
\end{equation}
оскільки чисельник обмежений, а знаменник $T$ зростає (або, для точно періодичного руху, $Q(T)=Q(0)$ й чисельник дорівнює нулю точно). Саме
ця властивість і дозволяє прибрати нижче доданки, які самі є повними похідними за часом.

Розкладаючи знаменник за малим параметром $r_i/r$,
\begin{equation*}
	\frac{1}{|\vect r-\vect r_i|} \approx \frac{1}{r} + \frac{\vect r\cdot\vect r_i}{r^3},
\end{equation*}
отримуємо
\begin{equation}\label{eq:VectorPotentialTwoTerms}
	\vect A(\vect r) = \frac{1}{cr}\sum_i q_i\vect v_i + \frac{1}{cr^3}\sum_i q_i\vect v_i(\vect r\cdot\vect r_i).
\end{equation}
Перший доданок є повною похідною за часом:
\begin{equation*}
	\sum_i q_i\vect v_i = \frac{d}{dt}\sum_i q_i\vect r_i,
\end{equation*}
причому $\sum_i q_i\vect r_i$ --- це електричний дипольний момент системи \eqref{eq:el_momentum}, обмежена величина. Тож, згідно
з~\eqref{eq:AveragedDerivativeZero}, цей член зникає після усереднення за часом --- саме той факт, який був анонсований у
§~\ref{sec:MagMoment}: аналога електричного заряду для струмів не існує навіть формально, оскільки єдиний природний кандидат на цю роль тотожно
зводиться до похідної дипольного моменту.

Для другого доданка позначимо $S \equiv \sum_i q_i\vect v_i(\vect r\cdot\vect r_i)$ і розглянемо допоміжну повну похідну
\begin{equation*}
	\frac{d}{dt}\sum_i q_i\vect r_i(\vect r\cdot\vect r_i)
	= \sum_i q_i\vect v_i(\vect r\cdot\vect r_i) + \sum_i q_i\vect r_i(\vect r\cdot\vect v_i)
	= S + \sum_i q_i\vect r_i(\vect r\cdot\vect v_i),
\end{equation*}
середнє від якої також дорівнює нулю (сума $\sum_i q_i\vect r_i(\vect r\cdot\vect r_i)$ обмежена), отже
\begin{equation}\label{eq:AuxRelation}
	\sum_i q_i\vect r_i(\vect r\cdot\vect v_i) = -S.
\end{equation}
Скориставшись тотожністю подвійного векторного добутку $\vect r\times(\vect v_i\times\vect r_i) = \vect v_i(\vect r\cdot\vect r_i) - \vect r_i(\vect r\cdot\vect v_i)$, підсумовуємо з вагами $q_i$ і враховуємо~\eqref{eq:AuxRelation}:
\begin{multline}
	\sum_i q_i\bigl[\vect r\times(\vect v_i\times\vect r_i)\bigr] = S - (-S) = 2S
	\Longrightarrow\\
	S = \frac12\,\vect r\times\sum_i q_i(\vect v_i\times\vect r_i) = -\vect r\times\left(\frac12\sum_i q_i(\vect r_i\times\vect v_i)\right).
\end{multline}
Сума в дужках тут --- це магнітний момент системи точкових зарядів~\eqref{eq:MagMomentDiscrete},
тому $S = -c\,(\vect r\times\vect p_m) = c\,(\vect p_m\times\vect r)$, і підстановка в~\eqref{eq:VectorPotentialTwoTerms} одразу дає:
\begin{equation}\label{eq:VectorPotentialFar}
	\tcbhighmath{
		\vect A(\vect r) = \frac{\vect p_m\times\vect r}{r^3}.
	}
\end{equation}
Магнітне поле далекої системи струмів знаходимо за формулою $\Bfield = \rot\vect A$, підставляючи~\eqref{eq:VectorPotentialFar}. Скористаємось тотожністю ротора векторного добутку
\begin{equation*}
	\rot(\vect a\times\vect w) = \vect a(\divg\vect w) - \vect w(\divg\vect a) + (\vect w\cdot\nabla)\vect a - (\vect a\cdot\nabla)\vect w,
\end{equation*}
застосованою до $\vect a=\vect p_m=\const$ і $\vect w = \vect r/r^3$. Оскільки $\vect p_m$ стала, $\divg\vect p_m=0$ і $(\vect w\cdot\nabla)\vect p_m=0$, тож
\begin{equation*}
	\Bfield = \rot\vect A = \vect p_m\,\divg\!\left(\frac{\vect r}{r^3}\right) - (\vect p_m\cdot\nabla)\frac{\vect r}{r^3}.
\end{equation*}
Перший доданок зникає, бо $\divg(\vect r/r^3)=0$ при $r\neq0$. У другому обчислюємо похідну по компонентах:
\begin{equation*}
	(\vect p_m\cdot\nabla)\frac{r_l}{r^3}
	= p_{m,k}\partial_k\!\left(\frac{r_l}{r^3}\right)
	= p_{m,k}\left(\frac{\delta_{kl}}{r^3} - \frac{3r_l r_k}{r^5}\right)
	= \frac{p_{m,l}}{r^3} - \frac{3r_l(\vect p_m\cdot\vect r)}{r^5},
\end{equation*}
тобто
\begin{equation*}
	(\vect p_m\cdot\nabla)\frac{\vect r}{r^3} = \frac{\vect p_m}{r^3} - \frac{3(\vect p_m\cdot\vect r)\vect r}{r^5}.
\end{equation*}
Підставляючи назад, отримуємо шуканий результат:
\begin{equation}\label{eq:DipoleFieldB}
	\tcbhighmath{
		\Bfield = \frac{3(\vect p_m\cdot\vect r)\vect r}{r^5} - \frac{\vect p_m}{r^3}.
	}
\end{equation}

\begin{Attention}
	Формула~\eqref{eq:DipoleFieldB} \emphx{за виглядом повністю збігається} з формулою для поля точкового електричного диполя, $\Efield =
	3(\vect p_e\cdot\vect r)\vect r/r^5 - \vect p_e/r^3$ (гл.~\ref{FieldInDielectrics}): достатньо замінити $\vect p_e \to \vect p_m$ і $\Efield
	\to \Bfield$. Це означає, що на достатньо великій відстані поле \emphx{будь-якої} системи струмів (колового витка, стрижневого магніту,
	навіть атома) виглядає так само, як поле точкового \emphx{магнітного диполя} --- незалежно від конкретного мікроскопічного розподілу цих
	струмів.
\end{Attention}

%% --------------------------------------------------------
\subsection*{Магнітний диполь і ефективні магнітні заряди}
%% --------------------------------------------------------

Оскільки формально поле магнітного диполя збігається з полем електричного, точковий магнітний момент можна (лише \emphx{формально}, поза
областю, зайнятою реальними струмами) розглядати як диполь, складений з \emphx{ефективних магнітних зарядів} --- північного ($N$) і південного
($S$), рис.~\ref{pic:DipoleCharges}:
\begin{equation*}
	\vect p_m = q_m\vect\ell,
\end{equation*}
де $q_m$ --- величина ефективного магнітного заряду, а $\vect\ell$ --- \enquote{плече} диполя.

\begin{SCfigure}[9][h!]\centering
	\localinput{dipole-equivalent-charges.tikz}
	\caption{Формальна модель магнітного диполя як пари ефективних магнітних зарядів
		$N$ і $S$\label{pic:DipoleCharges}}
\end{SCfigure}

\begin{Attention}
	Ця модель --- зручна \emphx{розрахункова фікція}, а не фізична реальність: реальних магнітних зарядів немає (теорема Гаусса для $\Bfield$,
	формула~\eqref{eq:GaussMagDiff}). Модель ефективних зарядів коректно відтворює поле \emphx{зовні} області, зайнятої струмами, але дає хибну
	картину всередині цієї області: лінії $\Bfield$ реального витка зі струмом --- \emphx{замкнені} і всередині витка напрямлені \emphx{в той
	самий бік}, що й зовні на осі; натомість лінії поля формальної моделі \enquote{стрижня} з ефективними зарядами на кінцях --- \emphx{розімкнені}
	і всередині йдуть від $N$ до $S$, тобто \emphx{назустріч} справжньому полю витка.
\end{Attention}


\begin{table}[h!]\small\centering
	\caption{Порівняння електричного і магнітного точкових диполів}
		\begin{tblr}{
		colspec={X[l, m]X[c,m]X[c,m]},
        row{odd} = {gray!10},
		row{1} = {c, m, fg=white, bg=themecolorlight, font=\bfseries},
		hlines,
		vlines
		}
		         & Електричний диполь                                                    & Магнітний диполь                                                     \\
		Потенціал & $\phi = \dfrac{\vect p_e\cdot\vect r}{r^3}$                          & $\vect A = \dfrac{\vect p_m\times\vect r}{r^3}$                      \\
		Поле      & $\Efield = \dfrac{3(\vect p_e\cdot\vect r)\vect r}{r^5} - \dfrac{\vect p_e}{r^3}$ & $\Bfield = \dfrac{3(\vect p_m\cdot\vect r)\vect r}{r^5} - \dfrac{\vect p_m}{r^3}$ \\
	\end{tblr}
\end{table}

%% --------------------------------------------------------
\section{Потенціальна енергія диполя та сила, що діє на диполь у магнітному полі}\hypertarget{magnetic-dipole-energy-force}{}
%% --------------------------------------------------------

%% --------------------------------------------------------
\subsection*{Момент сили, що діє на контур у магнітному полі}
%% --------------------------------------------------------

Розглянемо плоский виток зі струмом $I$ в \emphx{однорідному} зовнішньому полі $\Bfield$ (рис.~\ref{pic:TorqueLoop}). За означенням моменту сили і
формулою для сили Ампера на елементі витка~\eqref{eq:AmpereForceLin},
\begin{equation*}
	\vect M = \frac1c \oint\limits_L \vect r\times\left(I\, d\vect\ell\times\Bfield\right).
\end{equation*}

\begin{SCfigure}[2][h!]\centering
	\localinput{torque-loop-orientations.tikz}
	\caption{Виток зі струмом, довільно орієнтований в однорідному полі\label{pic:TorqueLoop}}
\end{SCfigure}

Розкриваємо подвійний векторний добуток за формулою \enquote{БАЦ мінус ЦАБ}:
\begin{equation*}
	\vect r\times(d\vect r\times\Bfield) = d\vect r\,(\vect r\cdot\Bfield) - \Bfield\,(\vect r\cdot d\vect r),
\end{equation*}
звідки
\begin{equation}\label{eq:TorqueTwoTerms}
	\vect M = \frac{I}{c}\oint\limits_L d\vect r\,(\vect r\cdot\Bfield) - \frac{I}{c}\Bfield\oint\limits_L(\vect r\cdot d\vect r).
\end{equation}
Другий доданок зникає одразу: підінтегральний вираз $\vect r\cdot d\vect r = d(r^2/2)$ --- повний диференціал, а інтеграл повного диференціала
по замкненому контуру дорівнює нулю:
\begin{equation}\label{eq:TotalDiffZero1}
	\oint\limits_L \vect r\cdot d\vect r = \oint\limits_L d\!\left(\frac{r^2}{2}\right) = 0.
\end{equation}

Залишається обчислити перший доданок, $\oint_L d\vect r\,(\vect r\cdot\Bfield)$. Розкладемо диадний вираз $dr_i\,r_j$ на симетричну й
антисиметричну за індексами $i,j$ частини:
\begin{equation*}
	dr_i\, r_j = \underbrace{\tfrac12(dr_i\,r_j + r_i\,dr_j)}_{\text{симетрична частина}} + \underbrace{\tfrac12(dr_i\,r_j - r_i\,dr_j)}_{\text{антисиметрична частина}}.
\end{equation*}
Симетрична частина є повним диференціалом добутку координат, $\tfrac12(dr_i r_j+r_i dr_j)=\tfrac12\,d(r_ir_j)$, і тому інтегрується в нуль по
замкненому контуру --- так само, як і в~\eqref{eq:TotalDiffZero1}:
\begin{equation*}
	\oint\limits_L \tfrac12\,d(r_i r_j)\,B_j = 0.
\end{equation*}
Антисиметрична частина виражається через векторний добуток: використовуючи тотожність $r_i\,dr_j - r_j\,dr_i = \varepsilon_{ijk}(\vect r\times
d\vect r)_k$, отримуємо
\begin{equation*}
	\tfrac12(dr_i r_j - r_i dr_j)\,B_j = \tfrac12\bigl[(\vect r\times d\vect r)\times\Bfield\bigr]_i,
\end{equation*}
тож
\begin{equation}\label{eq:AntisymmPart}
	\oint\limits_L d\vect r\,(\vect r\cdot\Bfield) = \tfrac12\left(\oint\limits_L\vect r\times d\vect r\right)\times\Bfield.
\end{equation}
Підставляючи~\eqref{eq:TotalDiffZero1} і~\eqref{eq:AntisymmPart} у~\eqref{eq:TorqueTwoTerms}, отримуємо:
\begin{equation}\label{eq:TorqueLoop}
	\tcbhighmath{
		\vect M = \left(\frac{I}{2c}\oint\limits_L \vect r\times d\vect r\right)\times\Bfield = \vect p_m\times\Bfield,
	}
\end{equation}
--- той самий результат, що й був покладений в означення вектора $\Bfield$ на початку розділу (формула~\eqref{eq:TorqueDipoleAnalogy}), але тепер
уже \emphx{виведений}, а не постульований, для довільного плоского витка довільної форми. Тут $\vect p_m = \frac{I}{2c}\oint_L\vect r\times d\vect
r = \frac{I}{c}\vect S$, де $\vect S=\frac12\oint_L\vect r\times d\vect r$ --- векторна площа контуру: це саме та формула магнітного моменту
лінійного струму~\eqref{eq:MagMomentLin}, яку ми означили ще на початку цього розділу (тут просто $d\vect r$ відіграє роль елемента контуру
$d\vect\ell$).

%% --------------------------------------------------------
\subsection*{Потенціальна енергія диполя в магнітному полі}
%% --------------------------------------------------------

Момент сил $\vect M = \vect p_m\times\Bfield$ прагне зорієнтувати магнітний момент витка вздовж поля. Якщо кут $\theta$ між $\vect p_m$ і $\Bfield$ збільшується на $d\theta$, момент сил виконує роботу $\delta A = -p_mB\sin\theta\, d\theta$, що дорівнює зменшенню потенціальної енергії, $\delta A = -dU$. Звідси $U = -p_mB\cos\theta$ (з точністю до сталої), тобто
\begin{equation}\label{eq:PotentialEnergyDipole}
	\tcbhighmath{
		U = -\vect p_m\cdot\Bfield.
	}
\end{equation}
Мінімум цієї енергії досягається при $\vect p_m\uparrow\uparrow \Bfield$ --- саме тому магнітний момент, якщо йому не заважають інші сили,
прагне орієнтуватись вздовж зовнішнього поля.

\begin{Attention}
	Формула~\eqref{eq:PotentialEnergyDipole}, разом із твердженням про нульову роботу магнітного поля~\eqref{eq:WorkMagnetic}, породжує
	позірний парадокс: якщо магнітний момент витка --- це реальний виток зі струмом, а магнітне поле, орієнтуючи цей виток, змінює його
	потенціальну енергію, --- то \emphx{за рахунок чого} відбувається ця зміна енергії, якщо магнітне поле принципово не може виконувати роботу?
	Відповідь --- зміна енергії витка при його повороті відбувається \emphx{за рахунок
		джерела струму}, яке підтримує в контурі струм $I$ незмінним, долаючи ЕРС індукції, що виникає під час повороту витка. Саме на цьому
	принципі --- перетворенні енергії джерела в механічну роботу обертання рамки за участі магнітного поля --- і ґрунтується
	принцип роботи \emphx{електричного двигуна}: ротор (рамка зі струмом), поміщений між полюсами магніту, під дією моменту сил $\vect
	M=\vect p_m\times\Bfield$ намагається зорієнтуватися вздовж поля, а комутатор своєчасно змінює напрям струму в обмотці, щоб забезпечити
	неперервне обертання в одному напрямку.
\end{Attention}

%%% --------------------------------------------------------
\subsection*{Сила, що діє на диполь у неоднорідному магнітному полі}
%%% --------------------------------------------------------


У \emphx{однорідному} полі, як бачимо з формули~\eqref{eq:TorqueLoop}, на диполь діє лише момент сил, що обертає його, але не діє сумарна сила
(сили Ампера, що діють на протилежні елементи витка, компенсують одна одну). У \emphx{неоднорідному} полі ця компенсація неповна, і на диполь діє
сумарна сила (рис.~\ref{pic:ForceDipole}):
\begin{equation*}
	\vect F = -\grad U = \grad(\vect p_m\cdot\Bfield).
\end{equation*}
Скориставшись векторною тотожністю $\grad(\vect a\cdot\vect b) = \vect b\times\rot\vect a + \vect a\times\rot\vect b + (\vect b\cdot\grad)\vect a
+ (\vect a\cdot\grad)\vect b$ і врахувавши, що магнітний момент $\vect p_m$ --- стала величина (тому $\rot\vect p_m \equiv 0$ і $(\Bfield\cdot\grad)\vect p_m = 0$), а в області, де
перебуває диполь, немає струмів провідності ($\rot\Bfield = 0$, оскільки $\vect j=0$), отримуємо остаточно
\begin{equation}\label{eq:ForceDipole}
	\tcbhighmath{
		\vect F = (\vect p_m\cdot\grad)\Bfield.
	}
\end{equation}

\begin{SCfigure}[2.5][h!]\centering
	\localinput{force-nonuniform-field.tikz}
	\caption{Магнітний диполь у неоднорідному полі втягується в область сильнішого поля\label{pic:ForceDipole}}
\end{SCfigure}

\begin{Attention}
	Формула~\eqref{eq:ForceDipole} означає, що магнітний диполь, зорієнтований уздовж поля, завжди \emphx{втягується} в область, де поле
	\emphx{сильніше} (там, де $B$ зростає), --- точнісінько так само, як електричний диполь у неоднорідному електричному полі. Саме на цьому
	принципі ґрунтується, наприклад, дія магніту на феромагнітні матеріали (цвяхи, залізні ошурки тощо): зовнішнє поле магніту наводить у них
	магнітний момент, зорієнтований уздовж поля, і цей момент після цього втягується в область найсильнішого поля --- впритул до полюса магніту.
\end{Attention}

%% --------------------------------------------------------
\section*{Підсумки}
\phantomsection
\addcontentsline{toc}{section}{Підсумки}
%% --------------------------------------------------------

Магнітне поле створюється \emphx{рухом} зарядів --- електричними струмами --- і діє лише на \emphx{рухомі} заряди. Це підтверджують і досліди Ерстеда, і досліди Роуленда та Ейхенвальда: заряджене тіло, що рухається, створює поле так само, як провідник зі струмом. Водночас поділ електромагнітного поля на електричну й магнітну складові залежить від системи відліку: заряд, що для одного спостерігача створює лише електричне поле, для іншого створює й магнітне. Тому $\Efield$ і $\Bfield$ --- це дві грані єдиного електромагнітного поля; повністю цей факт розкриває спеціальна теорія відносності (гл.~\ref{Relativity}).

Магнітних зарядів у природі не виявлено, тому вектор $\Bfield$ вводять через момент сил, що діє на магнітний диполь. Дію поля на заряди й струми описують силою Лоренца та силою Ампера. Принципово важливо, що магнітна складова сили Лоренца завжди перпендикулярна до швидкості й тому не виконує роботи: магнітне поле змінює напрям швидкості частинки, але не її кінетичну енергію.

Поле заданого розподілу струмів можна знайти двома шляхами: інтегруванням закону Біо--Савара--Лапласа з урахуванням принципу суперпозиції (зручно також користуватися вектор-потенціалом $\vect A$) або, для симетричних розподілів (прямий провідник, соленоїд, циліндр), за теоремою про циркуляцію. Диференціальні теореми $\divg\Bfield=0$ і $\rot\Bfield=4\pi\vect j/c$ визначають структуру магнітостатики: силові лінії поля замкнені, а саме поле є вихровим і породжується струмами. Порівняння з електростатикою (табл.~\ref{tab:ElMagCompare}) показує, що ця відмінність --- прямий наслідок відсутності магнітних зарядів.

На великій відстані від системи струмів поле визначається \emphx{магнітним моментом}: аналога електричного заряду для струмів не існує, тому дипольний член розкладу є головним. За виглядом поле магнітного диполя збігається з полем електричного. У зовнішньому полі на диполь діють момент сил, що прагне орієнтувати його вздовж поля, і, якщо поле неоднорідне, сила, що втягує диполь в область сильнішого поля. Цим пояснюють орієнтацію магнітної стрілки, дію магніту на залізні предмети й роботу електродвигуна.

Позірний парадокс --- поле не виконує роботи, а рамка в двигуні й феромагнетик поблизу магніту рухаються --- розв'язується так: механічна робота виконується за рахунок енергії джерела струму (або зовнішньої сили), а магнітне поле не є ні джерелом, ні сховищем цієї енергії для окремої частинки. Воно визначає напрям і величину сил, що діють на струми, і саме тому джерело струму здатне віддавати енергію рухові рамки. У наступних розділах ми вивчимо магнітне поле в речовині (гл.~\ref{MagFieldMedia}), явище електромагнітної індукції, де змінне магнітне поле породжує електричне (гл.~\ref{ElMagInduction}), і зведемо всі закони в систему рівнянь Максвелла (гл.~\ref{MaxwellEqns}).


%% --------------------------------------------------------
\section*{Задачі}
\phantomsection
\addcontentsline{toc}{section}{Задачі}
%% --------------------------------------------------------

\begin{problem}\label{prob:LoopAxisField}
	\emphx{Поле колового витка.} За законом Біо--Савара--Лапласа знайдіть індукцію магнітного поля на осі колового витка радіуса $R$ зі струмом $I$ на відстані $z$ від його центру. Чому дорівнює поле в центрі витка? Покажіть, що при $z\gg R$ результат збігається з полем точкового магнітного диполя \eqref{eq:DipoleFieldB} з моментом $p_m=\pi R^2I/c$.
%	\emphx{Відповідь:} $B=\dfrac{2\pi R^2I}{c\,(R^2+z^2)^{3/2}}$; у центрі $B=\dfrac{2\pi I}{cR}$. При $z\gg R$ маємо $B\approx\dfrac{2p_m}{z^3}$, що збігається з \eqref{eq:DipoleFieldB} на осі диполя ($\vect r\parallel\vect p_m$).
\end{problem}

\begin{problem}\label{prob:CylinderCavity}
	\emphx{Циліндр із порожниною.} Нескінченний циліндричний провідник радіуса $R$ зі сталою густиною струму $j$ (уздовж осі) має циліндричну порожнину, вісь якої паралельна осі провідника й зсунута відносно неї на відстань $d$ (порожнина повністю лежить усередині провідника). Знайдіть магнітне поле в порожнині. Підказка: подайте порожнину як накладання струмів $+\vect j$ і $-\vect j$.
	\emphx{Відповідь:} $\Bfield=\dfrac{2\pi}{c}\,[\vect j\times\vect a]$, де $\vect a$ --- вектор, проведений від осі провідника до осі порожнини.
%    Поле однорідне, $B=\dfrac{2\pi jd}{c}$, перпендикулярне до $\vect j$ і до лінії, що сполучає осі. (Усередині суцільного циліндра $\Bfield=\dfrac{2\pi}{c}[\vect j\times\vect r]$, де $\vect r$ --- відстань від осі.)
\end{problem}

\begin{problem}\label{prob:RotatingDiskMoment}
	\emphx{Магнітний момент і момент імпульсу.} (а)~Тонкий диск радіуса $R$ і маси $M$ рівномірно заряджений по поверхні зарядом $Q$ і обертається з кутовою швидкістю $\omega$ навколо осі, перпендикулярної до диска й проведеної через його центр. Знайдіть магнітний момент $p_m$, розбивши диск на кільця, і перевірте результат за формулою \eqref{eq:MagMomentViaInertia}. (б)~Електрон рухається по коловій орбіті з моментом імпульсу $L=\hbar$. Знайдіть його магнітний момент і взаємну орієнтацію $\vect p_m$ та $\vect L$ ($\hbar=1{,}05\cdot10^{-27}$~ерг$\cdot$с).
%	\emphx{Відповідь:} (а)~$p_m=\dfrac{QR^2\omega}{4c}$; за формулою \eqref{eq:MagMomentViaInertia} з $J=\tfrac12MR^2$ --- те саме. (б)~$|\vect p_m|=\dfrac{e\hbar}{2mc}\approx9{,}3\cdot10^{-21}$~ерг/Гс (магнетон Бора); оскільки заряд електрона від'ємний, $\vect p_m$ напрямлений \emphx{протилежно} до $\vect L$.
\end{problem}

\begin{problem}\label{prob:DipoleNearWire}
	\emphx{Магнітний диполь біля струму.} Малий виток з магнітним моментом $\vect p_m$ розташований на відстані $r$ від нескінченного прямого провідника зі струмом $I$; вектор $\vect p_m$ напрямлений уздовж поля $\Bfield$ провідника (по дотичній до кола навколо нього). Знайдіть силу, що діє на виток. Що зміниться, якщо $\vect p_m$ напрямити протилежно? Перевірте результат, розглянувши виток як квадрат зі стороною $a\ll r$ і струмом $I'$, у площині якого лежить провідник, і підрахувавши сили Ампера на його сторони, паралельні провіднику.
	\emphx{Відповідь:} $F=\dfrac{2p_mI}{cr^2}$, виток притягується до провідника (диполь втягується в область сильнішого поля \eqref{eq:ForceDipole}); за $\vect p_m\uparrow\downarrow\Bfield$ сила та сама за модулем, але це відштовхування.
    %Перевірка: $F\approx\dfrac{2II'a}{c^2}\Bigl(\dfrac1r-\dfrac1{r+a}\Bigr)\approx\dfrac{2II'a^2}{c^2r^2}$, що збігається з $\dfrac{2p_mI}{cr^2}$ при $p_m=I'a^2/c$. Сили на сторони, перпендикулярні до провідника, взаємно скорочуються.
\end{problem}

\begin{problem}\label{prob:RowlandDisk}
	\emphx{Дослід Роуленда--Ейхенвальда.} Тонкий диск радіуса $R = 10$~см рівномірно заряджений по поверхні зарядом $Q = 3\cdot10^{3}$~СГСЕ ($\approx 1$~мкКл) і обертається з частотою $100$~об/с. Знайдіть магнітне поле в центрі диска та порівняйте його з полем Землі ($\approx 0{,}5$~Гс). Як зміниться поле при зміні напряму обертання?
	\emphx{Відповідь:} $B = \dfrac{2\pi\sigma\omega R}{c} = \dfrac{2Q\omega}{cR} \approx 1{,}3\cdot10^{-5}$~Гс, тобто приблизно в $4\cdot10^{4}$ разів менше за земне; при зміні напряму обертання напрям $\Bfield$ змінюється на протилежний.
\end{problem}


\begin{Questions}
\begin{quizlist}
	\item У чому полягав дослід Ерстеда 1820 р.? Чому це відкриття стало відправною точкою об'єднання електрики й магнетизму в єдину теорію?
	\item Чим магнітне поле принципово відрізняється від електричного щодо дії на нерухомі заряди? Чому підхід до означення поля через силу на пробний заряд, застосовний для $\Efield$, не працює безпосередньо для магнітного поля?
	\item Виведіть означення вектора магнітної індукції $\Bfield$ через момент сил, що діє на магнітний диполь, $\vect M = \vect p_m\times\Bfield$. Чому $\Bfield$ називають індукцією, а не напруженістю, --- у чому історична причина цієї термінологічної асиметрії?
	\item Запишіть магнітну складову сили Лоренца та повну силу Лоренца. Виведіть з неї силу Ампера для елемента об'ємного $\vect j\,dV$ та лінійного $I\,d\vect\ell$ струму.
	\item Доведіть, що магнітна складова сили Лоренца не виконує роботи над зарядженою частинкою, $\delta A = 0$. Як цей факт узгоджується з тим, що електродвигун виконує реальну механічну роботу внаслідок дії магнітного поля на рамку зі струмом?
	\item Сформулюйте закон Біо-Савара-Лапласа для елемента струму. У чому полягає аналогія й відмінність цього закону від закону Кулона?
	\item Покажіть, що той самий рухомий заряд для одного спостерігача створює лише електричне поле, а для іншого --- і електричне, і магнітне. Що це означає для розуміння $\Efield$ і $\Bfield$ як граней єдиного електромагнітного поля?
	\item Що показали досліди Роуленда й Ейхенвальда? Виведіть поле в центрі зарядженого диска, що обертається, і поясніть, чому воно таке мале.
	\item Виведіть за законом Біо--Савара--Лапласа поле нескінченного прямого провідника $B=2I/(cr)$ та поле скінченного відрізка провідника через кути $\alpha_1$, $\alpha_2$.
	\item Що встановили досліди Ампера про взаємодію паралельних провідників зі струмом? Виведіть формулу сили взаємодії на одиницю довжини $F/\ell = 2I_1I_2/(c^2d)$ і поясніть, як вона історично визначала одиницю ампера.
	\item Виведіть вираз для вектор-потенціала $\vect A$ магнітного поля з закону Біо--Савара--Лапласа, показавши, що $\Bfield=\rot\vect A$. У чому калібрувальна свобода вибору $\vect A$ і що таке кулонівське калібрування $\divg\vect A=0$?
	\item Виведіть теорему Гаусса для магнітного поля $\divg\Bfield=0$ з тотожності $\divg\rot\vect A\equiv0$. Що це означає для силових ліній $\Bfield$ порівняно з лініями $\Efield$?
	\item Виведіть теорему про циркуляцію магнітного поля $\rot\Bfield = (4\pi/c)\vect j$ в диференціальній та $\oint\Bfield\cdot d\vect r = (4\pi/c)\sum I_i$ в інтегральній формах. Як визначають знак струму за правилом правого гвинта?
	\item Застосовуючи теорему про циркуляцію, знайдіть поле нескінченного прямого провідника, поле нескінченного соленоїда та поле всередині й поза нескінченним циліндричним провідником зі струмом.
	\item Порівняйте математичну структуру електро- і магнітостатики у вакуумі (потенціал, теорема Гаусса, теорема про циркуляцію). Чому магнітне поле вихрове, а електростатичне --- потенціальне, і як це пов'язано з відсутністю магнітних зарядів?
	\item Дайте означення магнітного моменту системи струмів через об'ємну та лінійну густину струму. Виведіть магнітний момент плоского колового витка $\vect p_m = IS\vect n/c$.
	\item Що таке гіромагнітне відношення? Виведіть його класичне значення $\gamma = Q/(2Mc)$ і поясніть, чому воно вдвічі порушується для електрона (аномальний $g$-фактор).
	\item Виведіть вираз для вектор-потенціалу та поля далекого магнітного диполя $\Bfield = 3(\vect p_m\cdot\vect r)\vect r/r^5 - \vect p_m/r^3$. Чому на великій відстані поле будь-якої системи струмів виглядає як поле точкового магнітного диполя?
	\item У чому полягає модель ефективних магнітних зарядів для опису поля диполя? Чому вона є лише зручною розрахунковою фікцією, коректною тільки поза областю реальних струмів?
	\item Виведіть момент сили $\vect M = \vect p_m\times\Bfield$, що діє на плоский виток зі струмом у однорідному полі, через силу Ампера на елементах контуру.
	\item Виведіть потенціальну енергію магнітного диполя в полі $U=-\vect p_m\cdot\Bfield$. Внаслідок чого насправді відбувається зміна цієї енергії, якщо магнітне поле принципово не виконує роботи?
	\item Виведіть формулу сили, що діє на магнітний диполь у неоднорідному полі, $\vect F = (\vect p_m\cdot\grad)\Bfield$. Чому диполь, зорієнтований уздовж поля, завжди втягується в область сильнішого поля, і як це пояснює притягання феромагнітних матеріалів до магніту?
\end{quizlist}
\end{Questions}